\documentclass[12pt]{article}
\usepackage[T1]{fontenc}
\usepackage[utf8]{inputenc}
\usepackage{lmodern}
\usepackage{textcomp}
\usepackage{enumitem}
\usepackage{geometry}
\geometry{margin=1in}
\begin{document}
\begin{center}
Answer Key - Cybersecurity - Graduate Level Assessment 18 \
\small (Answers are ordered exactly as in the question paper.)
\end{center}

\section*{Multiple Choice}
\begin{enumerate}[label=\arabic*.]
\item \textbf{Answer:} C
\\ \textit{Options:} A) IM - Trojans; B) Backdoor Trojans; C) SMS Trojan; D) Ransom Trojan
\\ \textit{Note:} An SMS Trojan is a type of malware that exploits mobile devices to send unauthorized text messages, often resulting in financial charges to the user without their consent.
\item \textbf{Answer:} C
\\ \textit{Options:} A) Sender Site; B) Site; C) Receiver site; D) Conferencing
\\ \textit{Note:} The correct answer is "Receiver site" because encryption ensures that only the intended recipient can access the original message, requiring decryption at the receiver's end to restore it to its readable form.
\item \textbf{Answer:} B
\\ \textit{Options:} A) Port, network, and services; B) Network, vulnerability, and port; C) Passive, active, and interactive; D) Server, client, and network
\\ \textit{Note:} The correct answer identifies three essential types of scanning in cybersecurity-network scanning for identifying active devices, vulnerability scanning for detecting security weaknesses, and port scanning for discovering open ports and services-each of which plays a crucial role in assessing and securing networks.
\item \textbf{Answer:} C
\\ \textit{Options:} A) Eavesdropping; B) Cross-site scripting; C) Authentication; D) SQL Injection
\\ \textit{Note:} Authentication is a fundamental security mechanism designed to verify the identity of users and is not an exploit, which refers to a method used to take advantage of vulnerabilities in a system.
\item \textbf{Answer:} B
\\ \textit{Options:} A) 160; B) 128; C) 64; D) 54
\\ \textit{Note:} The correct answer is 128 because MD 5 generates a fixed-length message digest of 128 bits, which is equivalent to 32 hexadecimal characters, regardless of the size of the input message.
\end{enumerate}

\section*{True / False}
\begin{enumerate}[label=\arabic*.]
\setcounter{enumi}{5}
\item \textbf{Answer:} False
\\ \textit{Options:} A) True; B) False
\\ \textit{Note:} Public key encryption is generally slower than symmetric key cryptography due to the complex mathematical operations involved, making it less efficient for tasks that require rapid data encryption and decryption.
\item \textbf{Answer:} False
\\ \textit{Options:} A) True; B) False
\\ \textit{Note:} WPA 2 was designed to enhance security on wireless networks primarily through robust encryption and authentication protocols, not specifically to provide an easy method for device connections.
\item \textbf{Answer:} False
\\ \textit{Options:} A) True; B) False
\\ \textit{Note:} The Cotton Road was actually a darknet market used for the illegal trade of drugs and other illicit goods, not a platform for legal goods and services.
\item \textbf{Answer:} True
\\ \textit{Options:} A) True; B) False
\\ \textit{Note:} IM-Trojans are indeed malicious programs specifically crafted to infiltrate instant messaging applications and capture users' login information, thereby compromising their security and privacy.
\item \textbf{Answer:} False
\\ \textit{Options:} A) True; B) False
\\ \textit{Note:} The statement is false because to verify a digital signature, the receiver uses the sender's public key, not their own private key, to confirm the authenticity of the signature and the sender's identity.
\end{enumerate}

\section*{Long Form Response}
\begin{enumerate}[label=\arabic*.]
\setcounter{enumi}{10}
\item \textbf{Answer:} Stealthiness is a critical aspect of penetration testing, as it directly impacts the effectiveness of the assessment by minimizing the likelihood of detection by security systems. To achieve this, penetration testers employ various strategies such as using low- and-slow techniques to avoid triggering intrusion detection systems (IDS) and conducting activities during off-peak hours to reduce visibility. Additionally, they may obfuscate their traffic and utilize encrypted communication methods to blend in with legitimate network activity. The implications of these strategies are significant; by remaining undetected, testers can more accurately simulate real-world attacks and identify vulnerabilities that organizations may not have noticed otherwise. Ultimately, the effectiveness of penetration testing hinges on the careful balance between thoroughness and stealthiness, ensuring that security measures are rigorously evaluated without prematurely alerting defenders.
\\ \textit{Note:} correct because it emphasizes the significance of stealthiness in penetration testing by outlining strategies like low-and-slow techniques, timing of activities, traffic obfuscation, and encryption, all of which are essential for minimizing detection and enhancing the effectiveness of security assessments.
\item \textbf{Answer:} SNMP enumeration plays a crucial role in cybersecurity by allowing attackers and security professionals to gather detailed information about network devices and configurations. Tools such as IP Network Browser facilitate this process by leveraging the Simple Network Management Protocol (SNMP) to query devices for valuable data, including device types, operating systems, and active services. The implications of using such tools in network security assessments are significant; while they can aid in identifying vulnerabilities and misconfigurations, they also raise ethical concerns regarding unauthorized access and potential exploitation of sensitive information. Thus, the responsible use of SNMP enumeration tools is essential to balance effective security assessments with adherence to legal and ethical standards in cybersecurity practices.
\\ \textit{Note:} correct because it accurately identifies SNMP enumeration as a method for gathering critical network information through specific tools, highlighting both its utility in security assessments and the potential risks involved in its use.
\end{enumerate}

\vfill
\noindent\textit{End of Answer Key}
\end{document}